%(BEGIN_QUESTION)
% Copyright 2008, Tony R. Kuphaldt, released under the Creative Commons Attribution License (v 1.0)
% This means you may do almost anything with this work of mine, so long as you give me proper credit

Sam klatrer opp et 130 fot høyt tårn med en lærebok som veier 8 pund (0.25 slug). Tony klatrer opp det samme tårnet med en lærebok som veier 5 pund (0.16 slug). Både Sam og Tony slipper bøkene sine fra toppen av tårnet nøyaktig samtidig. Se bort fra luftmotstand på bøkene under fritt fall, og beregn følgende:

\begin{itemize}
\item{} Arbeid utført av Sam ved løfting av læreboken =
\item{} Kinetisk energi til Sams lærebok rett før den treffer bakken =
\item{} Hastighet til Sams lærebok rett før den treffer bakken =
\end{itemize}

\begin{itemize}
\item{} Arbeid utført av Tony ved løfting av læreboken =
\item{} Kinetisk energi til Tonys lærebok rett før den treffer bakken =
\item{} Hastighet til Tonys lærebok rett før den treffer bakken =
\end{itemize}

\underbar{file i00430}
%(END_QUESTION)





%(BEGIN_ANSWER)

\begin{itemize}
\item{} Arbeid utført av Sam ved løfting av læreboken = 1040 ft-lb
\item{} Kinetisk energi til Sams lærebok rett før den treffer bakken = 1040 ft-lb
\item{} Hastighet til Sams lærebok rett før den treffer bakken = 91.2 ft/s
\end{itemize}

\begin{itemize}
\item{} Arbeid utført av Tony ved løfting av læreboken = 650 ft-lb
\item{} Kinetisk energi til Tonys lærebok rett før den treffer bakken = 650 ft-lb
\item{} Hastighet til Tonys lærebok rett før den treffer bakken = 91.2 ft/s
\end{itemize}

\vskip 10pt

Siden kinetisk energi rett før boken treffer bakken skal være lik potensiell energi ved slipp (forutsatt null energitap på grunn av luftmotstand), kan vi løse for $v$ ganske enkelt:

$$E_k = {1 \over 2}mv^2 \hbox{\hskip 100pt} E_p = mgh$$

$${1 \over 2}mv^2 = mgh$$

$${1 \over 2}v^2 = gh$$

$$v^2 = 2gh$$

$$v = \sqrt{2gh}$$

%(END_ANSWER)





%(BEGIN_NOTES)


%INDEX% Physics, energy, work, power: potential energy transformed to kinetic

%(END_NOTES)
